% ============ 15 上同调维数 (原书 pp.105-107) ============
\chapter{上同调维数}\label{sec:15}
\index{上同调维数}
\index{域的代数闭包}
\index{Tsen 定理}
\index{Tate 定理}

簇(或概形) $X$ 称为有上同调维数 $c$，如果 $c$ 是使对 $r > c$ 与 $X_{\mathrm{et}}$ 上的所有挠层 $F$，$H^{r}(X_{\mathrm{et}}, F) = 0$ 的最小整数(或 $\infty$).pro-有限群 $G$ 的上同调维数 $c$ 类似地定义：对 $r > c$ 与离散挠 $G$-模 $M$，$H^{r}(G, M) = 0$.对域 $K$，
\[
\operatorname{cd}(K) \overset{\mathrm{def}}{=} \operatorname{cd}(\operatorname{Spec} K) = \operatorname{cd}(\operatorname{Gal}(K^{\mathrm{sep}}/K)).
\]
例如，若 $K$ 可分闭，则 $\operatorname{cd}(K) = 0$；若 $K$ 有限，则 $\operatorname{cd}(K) = 1$(尽管 $H^{2}(\operatorname{Gal}(K^{\mathrm{al}}/K), \mathbb{Z}) = \mathbb{Q}/\mathbb{Z} \neq 0$)；且 $\operatorname{cd}(\mathbb{R}) = 1$(因为循环群的上同调是周期的).

代数拓扑的一个标准结果说，对实维数 $d$ 的流形 $M$，对 $r > d$，$H^{r}(M, F) = 0$；复分析的一个标准结果说，对复维数 $d$(从而实维数 $2d$)的 Stein 流形 $M$，对 $r > d$，$H^{r}(M, F) = 0$.这些论断的类比在平展拓扑中成立.

\begin{theorem}\label{thm:15.1}
对代数闭域 $k$ 上的簇 $X$，
\[
\operatorname{cd}(X) \leq 2\dim(X).
\]
若 $X$ 仿射，则
\[
\operatorname{cd}(X) \leq \dim(X).
\]
事实上，正如我们后面将看到的，当 $X$ 完备时，$H^{2\dim X}(X, \mathbb{Z}/n\mathbb{Z}) \neq 0$.整个证明中，$m = \dim X$.
\end{theorem}

我们说层 $F$ 的支集在维数 $d$ 中，如果 $F = \bigcup i_{Z*}F_Z$，其中并遍历 $X$ 的维数 $d$ 的不可约闭子簇 $Z$，$i_Z$ 是包含 $Z \hookrightarrow X$，$F_Z$ 是 $Z$ 上的层.对这样的层，
\[
H^{r}(X, F) = \varinjlim H^{r}(X, i_{Z*}F_Z) = \varinjlim H^{r}(Z, F_Z).
\]

\subsection*{域的上同调维数}

证明的出发点是域的如下结果(由 Grothendieck 猜想，由 Tate 证明).

\begin{theorem}\label{thm:15.2}
设 $K$ 是其子域 $k$ 上超越次数为 $d$ 的域.则
\[
\operatorname{cd}(K) \leq \operatorname{cd}(k) + d.
\]
\end{theorem}

\begin{proof}
Tsen 定理是 $d = 1$，$k$ 代数闭的情形.一般情形由此出发作归纳证明.见 Shatz 1972, p.119，定理 28.
\end{proof}

\subsection*{定理~\ref{thm:15.1} 第一个论断的证明}

定理对 $X$ 的维数作归纳证明.于是，可以假定若 $F$ 的支集在维数 $d < m$ 中，则
\[
H^{r}(X, F) = 0 \quad (r > 2d).
\]
可以假定 $X$ 连通.设 $g \colon \xi \hookrightarrow X$ 是 $X$ 的一般点.由伴随性对应于 $g^{*}F$ 上恒等映射的映射 $F \to g_{*}g^{*}F$ 在一般茎上诱导同构 $F_{\bar{\xi}} \to (g_{*}g^{*}F)_{\bar{\xi}}$，因此该映射的核与余核的每个截面都有真闭子概形上的支集.于是，由归纳假设，只需对形式为 $g_{*}F$ 的层证明定理.

\subsection*{子簇的局部环与严格局部环}

设 $X$ 是域 $k$ 上的不可约代数簇.对 $X$ 的任何不可约闭子簇 $Z$，可以联系环 $\mathcal{O}_{X, Z}$.用簇的语言，它是 $\varinjlim \Gamma(U, \mathcal{O}_X)$，其中 $U$ 遍历 $X$ 的满足 $U \cap Z \neq \varnothing$ 的开子集；用概形的语言，它是一般点处的局部环.它的剩余域是 $Z$ 上有理函数域 $k(Z)$，它在 $k$ 上的超越次数为 $\dim Z$.此外，$k(Z)$ 在 $k$ 上可分，即存在 $k(Z)$ 中代数无关的元素 $T_1, \ldots, T_m$，使得 $k(Z)$ 是 $k(T_1, \ldots, T_m)$ 的可分代数扩张(否则 $Z$ 就不是簇).

\begin{lemma}\label{lem:15.3}
设 $A$ 是含域 $k$ 的 Hensel 局部环.设 $K$ 是 $A$ 的剩余域，并假定 $K$ 在 $k$ 上可分.则 $A$ 含域 $L$，它被商映射 $\pi \colon A \to K$ 同构地映到 $K$ 上.
\end{lemma}

\begin{proof}
设 $T_1, \ldots, T_m$ 是 $A$ 中这样的元素：$\pi(T_1), \ldots, \pi(T_m)$ 在 $k$ 上代数无关，且 $K$ 是 $k(\pi(T_1), \ldots, \pi(T_m))$ 的可分代数扩张.设 $L$ 是含 $k(T_1, \ldots, T_m)$ 的 $A$ 的极大子域---由 Zorn 引理，这样的域存在.因为 $\mathfrak{m}_{A} \cap L = \{0\}$，$\pi$ 把 $L$ 同构地映到 $K$ 的子域 $\pi L$ 上.设 $\alpha \in K$.则 $\alpha$ 是多项式 $f(T) \in L[T]$ 的单根.因为 $A$ 是 Hensel 的，存在 $\alpha' \in A$，使得 $\pi(\alpha') = \alpha$ 且 $\alpha'$ 是 $f(T)$ 的根.现在 $L[\alpha']$ 是含 $L$ 的 $A$ 的子域.因为 $L$ 极大，$\alpha' \in L$，这表明 $\alpha \in \pi L$.
\end{proof}

\begin{lemma}\label{lem:15.4}
设 $Z$ 是簇 $X$ 的闭不可约子簇，$A$ 是 $\mathcal{O}_{X, Z}$ 的 Hensel 化.$L$ 的分式域 $L$ 含 $k(Z)$，且在 $k(Z)$ 上的超越次数为 $\dim X - \dim Z$.
\end{lemma}

\begin{proof}
前一个引理表明 $L \supset k(Z)$，因此只需计算它的超越次数.但 $L$ 是 $k(X)$ 的代数扩张，从而在 $k$ 上的超越次数为 $\dim X$，而 $k(Z)$ 在 $k$ 上的超越次数为 $\dim Z$.
\end{proof}

\begin{lemma}\label{lem:15.5}
层 $R^{s}g_{*}F$ 的支集在维数 $\leq m - s$ 中.
\end{lemma}

\begin{proof}
设 $Z$ 是 $X$ 的闭不可约子簇，$z = \operatorname{Spec} k(Z)$ 是它的一般点.$k(Z)^{\mathrm{sep}}$ 的可分闭包的选取确定 $X$ 的几何点 $\bar{z} \to z \hookrightarrow X$，而 $\bar{z}$ 处的严格局部环是 $\mathcal{O}^{\mathrm{h}}_{X, Z}$ 的最大非分歧扩张.于是，它的剩余域是 $k(Z)^{\mathrm{sep}}$，它的分式域 $K_{\bar{z}}$ 是 $k(X)$ 的代数扩张.与上面相同的论证表明它在 $k(Z)^{\mathrm{sep}}$ 上的超越次数为 $\dim X - \dim Z$.由 (\ref{ex:8.5})，
\[
(R^{s}g_{*}F)_{\bar{z}} = H^{s}(K_{\bar{z}}, F),
\]
由 (\ref{thm:15.2})，对 $s > m - \dim Z$，它为零.于是，正如所断言的，$(R^{s}g_{*}F)_{\bar{z}} \neq 0$ 蕴含 $s \leq m - \dim Z$.
\end{proof}

由归纳假设，发现
\[
H^{r}(X, R^{s}g_{*}F) = 0 \quad \text{当 } s \neq 0 \text{ 且 } r > 2(m - s) \text{ 时}.
\]
在谱序列
\[
H^{r}(X, R^{s}g_{*}F) \Rightarrow H^{r+s}(\xi, F)
\]
中，末项 $H^{r+s}(\xi, F) = 0$(当 $r + s > n$ 时，再由 15.2).由此推出，当 $r > 2m$ 时，$H^{r}(X, g_{*}F) = 0$.

\begin{question}\label{question:15.6}
由此还推出上面谱序列中的复形
\[
E^{0, n}_2 \to E^{2, n-1}_2 \to \cdots \to E^{2n, 0}_2
\]
必须正合.它是什么？
\end{question}

\subsection*{第二个论断的证明}

第二个论断的证明需要从 $m = 1$ 的情形(在 14.12 中证明)出发的复杂归纳与极限论证.它需要紧合基变换定理 \S 17.
